\section{Conflict.lean}

\code{Conflict.lean} defines when two hunks or two diffs conflict. It gives
both logical specifications in \code{Prop} and executable boolean checks in
\code{Bool}. It also proves that the boolean checks match the logical
specifications.

\subsection{Imports and Namespace}

This module depends on \code{Defns.lean} for \code{Range}, \code{Hunk}, and
\code{Diff}, and uses \code{Mathlib.Tactic} for the short symmetry and
reflection proofs. All declarations remain in the \code{VersionControl}
namespace.

\subsection{Logical Conflict Definitions}

\begin{definition}[identicalChange]
\code{[identicalChange]} has type \code{Hunk -> Hunk -> Prop}. It states that
two hunks have the same \code{range}, the same \code{before} lines, and the
same \code{after} lines. The \code{seq} and \code{op} fields are not part of
this equality test.
\end{definition}

\begin{definition}[pointInsertConflict]
\code{[pointInsertConflict]} has type \code{Hunk -> Hunk -> Prop}. It states
that both hunks are insert hunks, both insert at the same start line, and their
\code{after} lists are different.
\end{definition}

\begin{definition}[hunksConflict]
\code{[hunksConflict]} has type \code{Hunk -> Hunk -> Prop}. It states that two
hunks conflict when they are not identical and either their ranges overlap or
they have a point-insert conflict.
\end{definition}

\begin{definition}[diffsConflict]
\code{[diffsConflict]} has type \code{Diff -> Diff -> Prop}. It states that
two diffs conflict when there exists a hunk in the left diff and a hunk in the
right diff such that the two hunks conflict.
\end{definition}

\begin{definition}[compatible]
\code{[compatible]} has type \code{Diff -> Diff -> Prop}. It states that every
hunk in the left diff is non-conflicting with every hunk in the right diff.
\end{definition}

\subsection{Decidable Instances}

\begin{definition}[Decidable identicalChange]
\code{[Decidable identicalChange]} supplies a decision procedure for
\code{identicalChange left right}.
\end{definition}

\begin{definition}[Decidable Range.overlaps]
\code{[Decidable Range.overlaps]} supplies a decision procedure for
\code{left.overlaps right}.
\end{definition}

\begin{definition}[Decidable pointInsertConflict]
\code{[Decidable pointInsertConflict]} supplies a decision procedure for
\code{pointInsertConflict left right}.
\end{definition}

\subsection{Boolean Conflict Checks}

\begin{definition}[hunkConflictFlag]
\code{[hunkConflictFlag]} has type \code{Hunk -> Hunk -> Bool}. It is the
executable hunk-conflict check. It returns \code{true} exactly when the hunks
are not identical and either their ranges overlap or they have a point-insert
conflict.
\end{definition}

\begin{definition}[diffConflictFlag]
\code{[diffConflictFlag]} has type \code{Diff -> Diff -> Bool}. It is the
executable diff-conflict check. It returns \code{true} exactly when some pair
of hunks, one from each diff, has \code{hunkConflictFlag = true}.
\end{definition}

\begin{definition}[decide]
\code{[decide]} converts a decidable proposition into a boolean. This file uses
\code{decide} to compute \code{identicalChange}, \code{Range.overlaps}, and
\code{pointInsertConflict} inside \code{hunkConflictFlag}.
\end{definition}

\subsection{Hunk Theorems}

\begin{theorem}[identicalChange\_comm]
\code{[identicalChange_comm]} states that \code{identicalChange} is symmetric:
\code{identicalChange left right <-> identicalChange right left}.
\end{theorem}

\begin{theorem}[not\_hunksConflict\_of\_identicalChange]
\code{[not_hunksConflict_of_identicalChange]} states that identical hunks do
not conflict.
\end{theorem}

\begin{theorem}[hunksConflict\_comm]
\code{[hunksConflict_comm]} states that hunk conflict is symmetric:
\code{hunksConflict left right <-> hunksConflict right left}.
\end{theorem}

\begin{theorem}[identical\_hunks\_not\_conflict]
\code{[identical_hunks_not_conflict]} states that a hunk does not conflict with
itself.
\end{theorem}

\begin{theorem}[separated\_hunks\_not\_conflict]
\code{[separated_hunks_not_conflict]} states that two hunks do not conflict if
their ranges are separated and there is no point-insert conflict.
\end{theorem}

\begin{theorem}[same\_point\_equal\_insert\_not\_conflict]
\code{[same_point_equal_insert_not_conflict]} states that two inserts at the
same point with the same inserted text do not conflict, assuming both
\code{before} lists are empty.
\end{theorem}

\subsection{Boolean Reflection Theorems}

\begin{theorem}[hunkConflictFlag\_eq\_true\_iff]
\code{[hunkConflictFlag_eq_true_iff]} states that
\code{hunkConflictFlag left right = true} is equivalent to
\code{hunksConflict left right}.
\end{theorem}

\begin{theorem}[hunkConflictFlag\_eq\_false\_iff]
\code{[hunkConflictFlag_eq_false_iff]} states that
\code{hunkConflictFlag left right = false} is equivalent to
\code{not hunksConflict left right}.
\end{theorem}

\begin{theorem}[identical\_hunks\_not\_flagged]
\code{[identical_hunks_not_flagged]} states that
\code{hunkConflictFlag hunk hunk = false}.
\end{theorem}

\begin{theorem}[same\_point\_equal\_insert\_not\_flagged]
\code{[same_point_equal_insert_not_flagged]} is the boolean version of
\code{same_point_equal_insert_not_conflict}. It states that equal same-point
insert hunks produce \code{hunkConflictFlag = false}.
\end{theorem}

\subsection{Diff Theorems}

\begin{theorem}[diffConflictFlag\_eq\_true\_iff]
\code{[diffConflictFlag_eq_true_iff]} states that
\code{diffConflictFlag left right = true} is equivalent to
\code{diffsConflict left right}.
\end{theorem}

\begin{theorem}[diffConflictFlag\_eq\_false\_iff]
\code{[diffConflictFlag_eq_false_iff]} states that
\code{diffConflictFlag left right = false} is equivalent to
\code{not diffsConflict left right}.
\end{theorem}

\begin{theorem}[flagged\_implies\_real\_conflict]
\code{[flagged_implies_real_conflict]} states that if the executable diff flag
returns \code{true}, then the logical proposition \code{diffsConflict} holds.
\end{theorem}

\begin{theorem}[real\_conflict\_implies\_flagged]
\code{[real_conflict_implies_flagged]} states that if \code{diffsConflict}
holds, then the executable diff flag returns \code{true}.
\end{theorem}

\begin{theorem}[diffsConflict\_comm]
\code{[diffsConflict_comm]} states that diff conflict is symmetric:
\code{diffsConflict left right <-> diffsConflict right left}.
\end{theorem}

\begin{theorem}[compatible\_iff\_not\_diffsConflict]
\code{[compatible_iff_not_diffsConflict]} states that \code{compatible left right}
is equivalent to \code{not diffsConflict left right}.
\end{theorem}

\begin{theorem}[compatible\_empty\_left]
\code{[compatible_empty_left]} states that an empty diff on the left is
compatible with any diff.
\end{theorem}

\begin{theorem}[compatible\_empty\_right]
\code{[compatible_empty_right]} states that an empty diff on the right is
compatible with any diff.
\end{theorem}

\begin{theorem}[compatible\_comm]
\code{[compatible_comm]} states that compatibility is symmetric:
\code{compatible left right <-> compatible right left}.
\end{theorem}

\subsection{Execution Pipeline}

The conflict-checking pipeline is:

\begin{enumerate}
  \item \code{hunkConflictFlag} computes conflict for two hunks.
  \item \code{diffConflictFlag} computes whether any hunk pair conflicts.
  \item \code{hunkConflictFlag_eq_true_iff} connects the hunk boolean check to
    \code{hunksConflict}.
  \item \code{diffConflictFlag_eq_true_iff} connects the diff boolean check to
    \code{diffsConflict}.
  \item \code{compatible_iff_not_diffsConflict} states compatibility as the
    absence of diff conflict.
\end{enumerate}
